\section{Theorie}

\subsection{Säure-Base-Theorie}

Das grundlegende Modell zur Beschreibung von Protolysen ist die Säure-Base-Theorie nach Brønsted. Sie definiert Säuren als Protonendonatoren und Basen als Protonenakzeptoren. Nach dieser Theorie können protolytische Reaktionen sowohl in der flüssigen als auch in der Gasphase ablaufen.

\subsection{Autoprotolyse des Wassers und pH-Wert}

Anhand der Autoprotolyse des Wassers lassen sich sowohl die Eigenschaften von Ampholyten als auch die Prinzipien protolytischer Reaktionen veranschaulichen. Dabei läuft die in \eqref{eq:protolyse_wasser} dargestellte Reaktion ab. \cite{skript}
\begin{equation}
    \ce{2H2O_{\mathrm{(l)}} <=> H3O+ _{\mathrm{(aq)}} + OH- _{\mathrm{(aq)}}}
    \label{eq:protolyse_wasser}
\end{equation}
Hierbei reagiert ein Wassermolekül als Protonenakzeptor und ein zweites als Protonendonator. Bei der Rückreaktion fungiert wiederum das entstehende Oxonium-Ion als Protonendonator und das Hydroxid-Ion als Protonenakzeptor \cite{skript}.

Für diese Gleichgewichtsreaktion \eqref{eq:protolyse_wasser} lässt sich nun das Massenwirkungsgesetz \eqref{eq:mwg_wasser} aufstellen.
\begin{equation}
    K = \frac{[\ce{H3O+}][\ce{OH-}]}{[\ce{H2O}]^2}
    \label{eq:mwg_wasser}
\end{equation}

%%%%%%%%%%%%%%%%%%%%%%%%%%%%%%%%%%
%Die Konzentration des Wassers ist dabei näherungsweise als konstant zu betrachten. Damit ergibt sich aus dem Massenwirkungsgesetz \eqref{eq:mwg_wasser} der $K_{\mathrm{w}}$-Wert \eqref{eq:KW-Wert} wie folgt.
%\begin{equation}
%    K_{\mathrm{w}} = [\ce{H+}][\ce{OH-}]
%    \label{eq:KW-Wert}
%\end{equation}
%%%%%%%%%%%%%%%%%%%%%%%%%%%%%%%%%%

Da die Konzentration des Wassers in verdünnten wässrigen Lösungen etwa \SI{55,5}{\frac{\mole}{l}} beträgt und folglich nahezu konstant ist, kann sie mit der Gleichgewichtskonstanten $K$ zu einer neuen Konstante zusammengefasst werden. Dadurch wird das Ionenprodukt des Wassers $K_{\mathrm{w}}$ definiert \eqref{eq:KW-Wert}. \cite{skript}

\begin{equation}
    K_{\mathrm{w}} = K \cdot [\ce{H2O}]^2 = [\ce{H3O+}] \cdot [\ce{OH-}] %kürzen?
    \label{eq:KW-Wert}
\end{equation}

Bei \SI{25}{\celsius} beträgt das Ionenprodukt des Wassers $K_\mathrm{W} = \SI{e-14}{  \frac{\mole\squared}{\liter\squared}}$. Da solche Ionenkonzentrationen mehrere Größenordnungen unter der Stoffmengenkonzentration des Wassers liegen, verwendet man als Maß für den sauren oder basischen Charakter einer Lösung den pH-Wert \eqref{eq:ph_definition}. \cite{skript}

\begin{equation}
    \mathrm{pH} = -\log [\ce{H3O+}]
    \label{eq:ph_definition}
\end{equation}

% pKw?!!!!!!!!!!!!!!!!!!!!!!!!!!!!!!!!!

%%%%%%%%%%%%%%%%%%%%%%%%%%%%%%%%%%%%%%%%%%%%%%%%%%%%%%%%%%%%%%%%%%%%%%%%%%%%%%%
%\subsection{Autoprotolyse des Wassers und pH-Wert}
%Wasser ist amphoter, das heißt es kann sowohl als Säure als auch als Base reagieren, und es kommt zur Autoprotolyse\cite{skript}:
%\begin{equation}
%    \ce{2H2O <=> H3O+ + OH-}
%    \label{eq:autoprotolyse}
%\end{equation}
%
%
%Das Ionenprodukt des Wassers lässt sich direkt aus dem Massenwirkungsgesetz ableiten, da die Konzentration des Wassers als konstant angenommen wird, und deswegen gilt:
%\begin{equation}
%    K_\mathrm{W} = [\ce{H3O+}] \cdot [\ce{OH-}]
%    \label{eq:kw_definition}
%\end{equation}
%
%
%Bei \SI{25}{\celsius} gilt $K_\mathrm{W} = \SI{1e-14}{\mole\squared\per\liter\squared}$.\\
%Für neutrales Wasser folgt:
%\begin{equation}
%    [\ce{H3O+}] = [\ce{OH-}] = 10^{-7}\,\si{\mole\per\liter}
%    \label{eq:neutralwasser}
%\end{equation}
%
%Der pH-Wert ist definiert als der negative dekadische Logarithmus der Oxoniumionenkonzentration \cite{skript}:
%\begin{equation}
%    \mathrm{pH} = -\log [\ce{H3O+}]
%    \label{eq:ph_definition}
%\end{equation}
%
%Es gilt außerdem der Zusammenhang zwischen dem pH-Wert und dem pOH-Wert \cite{skript}:
%\begin{equation}
%    \mathrm{p}K_\mathrm{W} = \mathrm{pH} + \mathrm{pOH}
%    \label{eq:pkw_beziehung}
%\end{equation}
%%%%%%%%%%%%%%%%%%%%%%%%%%%%%%%%%%%%%%%%%%%%%%%%%%%%%%%%%%%%%%%%%%%%%%%%%


%%%%%%%%%%%%%%%%%%%%%%%%%%%%%%%%%% - temp felix version
\subsection{Die Autoprotolyse des Wassers -temp felix k}

Die Autoprotolyse des Ampholyten Wasser ist ein anschauliches Beispiel
sowohl für die Eigenschaften eines Ampholyten als auch für Protolysen
insgesamt. Dabei findet Reaktion \eqref{eq:Autoprotolyse} statt\cite{skript}.
\begin{equation}
    \text{H}_{2}\text{O}_{(\text{l})}+\text{H}_{2}\text{O}_{(\text{l})}\rightleftharpoons\text{H}_{3}\text{O}^{+}_{(\text{aq})}+\text{OH}^{-}_{\text{(aq)}}\label{eq:Autoprotolyse}
\end{equation}

Hierbei reagiert das ein Wasserstoffmolekül als Protonenakzeptor und
ein anderes Molekül als Protonendonator. Bei der Rückreaktion reagiert
wiederum das Oxonium-Ion als Protonendonator und das Hydroxid-Ion
als Protonenakzeptor.\cite{skript}

Für diese Reaktion lässt sich nun das Massenwirkungsgesetz (MWG) \eqref{eq:mwg}
aufstellen.
\begin{equation}
    K=\frac{[\text{H}_{3}\text{O}^{+}][\text{OH}^{-}]}{[\text{H}_{2}\text{O}]^{2}}\label{eq:mwg}
\end{equation}
Aus diesem lässt sich dann durch Multiplikation mit der Konzentration
des Wassers, die aufgrund der Tatsache, dass nur sehr wenige Wassermoleküle
dissoziieren, nahezu konstant bleibt, folgern \cite{skript}:
\begin{equation}
    K_{\text{w}}=[\text{H}_{3}\text{O}^{+}][\text{OH}^{-}]\label{eq:Kw}
\end{equation}
Bei $T=25{^{\circ}}\text{C}$ ist das Ionenprodukt des Wassers $K_{\text{w}}=10^{-14}\,\frac{\text{mol}^{2}}{\text{l}^{2}}$.
Da in der Säure-Base-Chemie auftretende Zahlenwerte häufig sehr klein
sind, werden sie meistens stattdessen als p-Werte angegeben, also
als negativer dekadischer Logarithmus. Daraus folgt, dass $\text{p}K_{\text{W}}=14$.
Genauso gibt es entsprechend den pH-Wert \eqref{eq:ph}, der die Konzentration
an Oxoniumionen angibt, welcher bei Standardbedingungen bei 7 liegt.
\begin{equation}
    \text{pH}=-\log([\text{H}^{+}])\label{eq:ph}
\end{equation}
%%%%%%%%%%%%%%%%%%%%%%%%%%%%%%%%%%

\subsection{Säure-Base-Theorie nach Brønsted}

Die Brønsted-Theorie besagt, dass die Zugabe einer Säure zu einer wässrigen Lösung die Konzentration von Oxoniumionen erhöht, während die Zugabe einer Base diese senkt. Dabei kennzeichnet ein pH-Wert unter 7 eine saure, ein pH-Wert über 7 eine basische Lösung. Aus Gleichung Nr.\eqref{eq:Kw} und $K_{\text{w}}$ folgt, dass die Summe des pH-Wert und pOH-Wert bei Normalbedingungen immer 14 ist, wodurch sich der $\text{pH}$-Wert aus dem $\text{pOH}$-Wert und umgekehrt berechnen lassen. Zudem ist nach der Brønsted-Theorie jeder Säure und Base ein $K_{\text{S}}$-, bzw. $\text{p}K_{\text{S}}$-Wert\cite{skript}, welcher die Gleichgewichtskonstante $K$ des Reaktionsgleichgewichts zwischen einer Säure und deren konjugierter Base darstellt, zugeordnet. Dieser Wert wird aus der Gleichgewichtskonstante $K$ durch Multiplikation mit der Konzentration des Wasser berechnet. \newline
Aus $K_{\text{w}}$ und aus \eqref{eq:Kw}
folgt, dass die Summe aus pH-Wert und pOH-Wert unter Normalbedingungen
immer 14 ist. Damit lässt sich aus dem $\text{pH}$-Wert der $\text{pOH}$-Wert
und umgekehrt berechnen. Nach der Brønsted-Theorie ist außerdem jeder
Säure und jeder Base ein Wert zugeordnet, der die Gleichgewichtskonstante
$K$ der Reaktionsgleichgewichts zwischen Säure und konjugierter Base
repräsentiert. Dieser Wert ist der $K_{\text{S}}$-, bzw. $\text{p}K_{\text{S}}$-Wert
\cite{skript}. Dieser wird aus $K$ durch Multiplikation mit der
Konzentration des Wassers berechnet.
\begin{equation}
    K_{\text{S}}=\frac{[\text{H}_{3}\text{O}^{+}][\text{A}^{-}]}{[\text{HA}]}
\end{equation}
Daneben gibt es für Basen den $\text{p}K_{\text{B}}$-Wert, der mit
dem $\text{p}K_{\text{S}}$-Wert der konjugierten Säure gemäß Gleichung
\eqref{eq:pkbpks} verknüpft ist.
\begin{equation}
    \text{p}K_{\text{S}}+\text{p}K_{\text{B}}=14\label{eq:pkbpks}
\end{equation}

Bei starken Säuren, die nahezu bis vollständig dissoziieren, ergibt
sich der pH-Wert dann gemäß Gleichung \eqref{eq:phs}, denn hier wird
genähert, dass die initiale Konzentration der hinzugegebenen Säure
gleich der Konzentration Oxonium-Ionen nach Dissoziiation.
\begin{equation}
    \text{pH}=-\log(c_{0})\label{eq:phs}
\end{equation}
Der pH-Wert einer schwachen Säure ist \eqref{eq:phw} \cite{zsfg_maier}.
Die hier genutzte Vereinfachung ist, dass bei einem hohen $\text{p}K_{\text{S}}$-Wert
die Initialkonzentration der zugegebenen Säure und die spätere nicht
dissoziierte Konzentration nahezu identisch sind.
\begin{equation}
    \text{pH}=\frac{1}{2}(\text{p}K_{\text{S}}-\log(c_{0}))\label{eq:phw}
\end{equation}


\subsection{Säure- und Basenstärke}

Die Säurestärke beschreibt wie gut eine Säure Protonen abgeben kann, und wird durch die Säurekonstante beschrieben. Die Säurekonstante lässt sich aus dem Massenwirkungsgesetz für die Reaktion \ref{eq:brønsted_allgemein} ableiten, in dem die Konzentration des Wassers als konstant angenommen wird \cite{skript}:
\begin{equation}
    K_\mathrm{S} =
    \frac{[\ce{H3O+}][\ce{A^-}]}{[\ce{HA}]}
    \label{eq:ks_definition}
\end{equation}

Der p$K$S-Wert ist definiert als der negative dekadische Logarithmus der Säurekonstante \cite{skript}:
\begin{equation}
    \mathrm{p}K_\mathrm{S} = -\log K_\mathrm{S}
    \label{eq:pks_definition}
\end{equation}

Für schwache Säuren gilt mithilfe der Näherung, dass die Konzentration der Säure vor und nach der Reaktion nahezu gleich ist\cite{skript}:
\begin{equation}
    [\ce{H3O+}] = \sqrt{K_\mathrm{S} \cdot c_0}
    \label{eq:schwache_saeure_nahe}
\end{equation}

\subsection{Protolysegleichgewichte}

Für konjugierte Säure-Base-Paare gilt, die Beziehung lässt sich wie folgt herleiten \cite{atkins_kurzlehrbuch_2020}:
\begin{equation}
    K_\mathrm{S} \cdot K_\mathrm{B} = \frac{[\ce{H3O+}][\ce{A^-}]}{[\ce{HA}]} \cdot \frac{[\ce{OH-}][\ce{HA}]}{[\ce{A^-}]} = [\ce{H3O+}] \cdot [\ce{OH-}] = K_\mathrm{W}
    \label{eq:ks_kb_herleitung}
\end{equation}
\begin{equation}
    K_\mathrm{S} \cdot K_\mathrm{B} = K_\mathrm{W}
    \label{eq:ks_kb_beziehung}
\end{equation}

\subsection{Puffersysteme}

Ein Puffersystem besteht aus einem konjugierten Säure-Base-Paar, dabei kann die Säure Basen abfangen und die Base Säuren abfangen, dadurch bleibt der pH-Wert stabil. Der pH-Wert eines Puffers ergibt sich aus der Henderson-Hasselbalch-Gleichung \cite{skript}:

\begin{equation}
    \mathrm{pH} = \mathrm{p}K_\mathrm{S}
    + \log \left(
    \frac{[\ce{A^-}]}{[\ce{HA}]}
    \right)
    \label{eq:henderson_hasselbalch}
\end{equation}

Für gleiche Konzentrationen gilt:

\begin{equation}
    [\ce{A^-}] = [\ce{HA}] \Rightarrow \mathrm{pH} = \mathrm{p}K_\mathrm{S}
    \label{eq:puffer_halb}
\end{equation}

\subsection{Indikatoren}

Säure-Base-Indikatoren sind schwache Säuren oder Basen, deren protonierte und deprotonierte Form unterschiedliche Absorptionsspektren besitzen. Der Farbumschlag beruht auf einem pH-abhängigen Gleichgewicht zwischen diesen beiden Formen \cite{skript}.

Je nach pH-Wert der Lösung verschiebt sich dieses Gleichgewicht in Richtung der protonierten oder deprotonierten Form, wodurch der charakteristische Farbwechsel entsteht. Der Umschlagsbereich entspricht dabei dem pH-Bereich, in dem beide Formen in vergleichbaren Konzentrationen vorliegen. Indikatoren stehen im Gleichgewicht \cite{mortimer_muller}:

\begin{equation}
    \ce{HIn <=> H+ + In^-}
    \label{eq:indikator_gleichgewicht}
\end{equation}

\newpage
